\documentclass{template}


\begin{document}

\section*{Seminar Proposals}

\subsection*{Optimizing Elliptic Curve Arithmetic}

\begin{itemize}
    \item Elliptic Curves and their importance to Cryptography (Definition of a
    group $E(\mathbb{F}_p)$, Discrete Log and Diffie-Hellman assumptions). \textcolor{gray}{See
    ``A Graduate Course in Applied Cryptography'', Section 15,
    \url{https://crypto.stanford.edu/~dabo/cryptobook/BonehShoup_0_6.pdf}.}
    \item Projective Elliptic Curve equation (in homogeneous and Jacobian
    forms), addition law in projective space, why it optimizes classical
    addition law in affine space. \textcolor{gray}{See ``Complete addition
    formulas for prime order elliptic curves'',
    \url{https://eprint.iacr.org/2015/1060}}.
    \item Scalar multiplication using double-and-add and non-adjacent form
    algorithms. Multiple point multiplication (MSM) and its application to
    polynomial commitment schemes (KZG).
    \item Optimizing scalar multiplication using efficiently computable
    endomorphisms. GLV Scalar Multiplication. \textcolor{gray}{See ``Guide to
    Elliptic Curve Cryptography'', Section 3.5.}
\end{itemize}

\subsection*{Fast Fourier Transform over Finite Fields}

\begin{itemize}
    \item A short remark on structure of $\mathbb{F}_q^{\times}$ and primitive
    $N$-th roots of unity in $\mathbb{F}_q$.
    \item Lagrange and Barycentric Interpolation formulas. Using $N$-th roots of
    unity subgroup for reducing complexity assymptotic. 
    \item Fast Fourier Transform over FFT-friendly finite fields.
    \item Elliptic Curve Fast Fourier Transform (ECFFT). See ``ECFFT: Fast
    Polynomial Algorithms over all Finite Fields''
    \url{https://arxiv.org/pdf/2107.08473}.
\end{itemize}

\textit{Note:} The last point of the lecture requires most preparation. The
first three are already presented in my book. One might see a simplified
presentation on these topics here:
\url{https://zkdl-camp.github.io/files/slides/13-ntt.pdf}.

\subsection*{Introduction to Zero-Knowledge Proofs}

\begin{itemize}
    \item What are zero-knowledge proofs and their applications. Interactive
    Probabilistic Proofs (IP). Soundness, Argument of Knowledge, Zero-Knowledge
    properties. Fiat-Shamir Heuristic.
    \item $\Sigma$-protocols as an example of classical IP.
    \item A glimpse into pairing-based zero-knowledge succinct non-interactive
    argument of knowledge. Schwartz-Zippel Lemma, KZG polynomial commitment
    scheme. High-level overview of Groth16: algebraic arithmetization and
    quadratic arithmetic programs.
\end{itemize}

\textit{Note:} To avoid using pairings without knowing their algebraic nature,
instead of the last point, we might consider zk-SNARKs based on multilinear
extensions (Sum-Check and GKR protocols). See
\url{https://zkdl-camp.github.io/files/slides/15-sumcheck.pdf}.

\end{document}